\scanpage{355}
% top centre
\mspage{271}
% lines 1–4
\uline{Proposition} 4.11.5. Soit $\unsure{\underline{\Phi}|\Phi}$ un ensemble de \deleted{petites} catégories \inserted{$I$} équivalentes à des \deleted{petites} catégories \inserted{finies} et ``rigides'' i.e. pour tt $i\in\operatorname{Ob} I$, \deleted{\unsure{tout}} $\operatorname{Hom}_I(i,i)=\{\mathrm{id}_i\}$.
% upper left beside proposition
% slanted
\margin{? Démonstration pas achevée}
% lines 4–12
\deleted{Soit $C$} \unsure{Soit|Soient} $C$ une catégorie équivalente à une petite catégorie, $\pi$ un cardinal infini, $M=\operatorname{Ind}_\pi C$. Pour que $M$ soit stable par $\varinjlim$ (resp. par $\varprojlim$) de type $\Phi$, il \uline{suffit} que $C$ le soit. Dans le cas \unsure{non respé|non resp}, si $C$ est karoubienne, cette condition est aussi nécessaire.

% lines 13–16
Cette dernière assertion est contenue dans l'assertion similaire de 4.11.1. Pour la première assertion, elle résulte de la façon habituelle du

% lines 17–20
\uline{Lemme} 4.11.6. Soit $J$ une catégorie équivalente à une \deleted{petite} catégorie finie, $J$ rigide, \deleted{Alors} $C$ une catégorie équiv. à une petite catégorie, $\pi$ un cardinal infini.
% lower left beside lemma
% slanted
\margin{? Démonstration \unsure{non|pas} \unsure{explicitée}, c'est peut-être faux !}
% lines 21–24
Alors \uline{tout} foncteur \deleted{\unsure{$I$|$J$}} $J\longrightarrow M\overset{\mathrm{déf}}{=}\operatorname{Ind}_\pi C$ admet une représentation indicielle $J\times I\longrightarrow C$, avec $I$ grde devant $\pi$, en d'autres termes le
% archive sheet number: 354
